\documentclass[11pt]{article}

\usepackage[a4paper,margin=29mm]{geometry}
\usepackage[T1]{fontenc}
\usepackage{lmodern}
\usepackage{microtype}
\usepackage{amsmath,amssymb,amsthm,mathtools}
\usepackage{booktabs}
\usepackage{enumitem}
\usepackage[hidelinks]{hyperref}
\setlength{\emergencystretch}{3em}

\newtheorem{theorem}{Theorem}[section]
\newtheorem{proposition}[theorem]{Proposition}
\newtheorem{corollary}[theorem]{Corollary}
\newtheorem{lemma}[theorem]{Lemma}
\theoremstyle{remark}
\newtheorem{remark}[theorem]{Remark}

\newcommand{\Sym}{\operatorname{Sym}}
\newcommand{\per}{\operatorname{per}}

\title{Orientation and exchange statistics in a Golden six-mode transfer}
\author{Tavis Rudd}
\date{August 2026}

\begin{document}
\maketitle

\begin{abstract}
Exchange statistics determine which information a many-particle
interferometer can retain about its one-particle transfer.  We study this
distinction in an exact six-mode transfer constructed from the two
three-dimensional eigenspaces of a real symmetric conference matrix.  For a
diagonal path control \(D_x\), the postselected block

\[
 K_T(x)=Q_{T,-}^{\mathsf T}D_xQ_{T,+}
\]

has filled three-fermion amplitude \(\det K_T(x)\), proportional to a signed
Joubert coordinate on the Segre cubic.  We show that this orientation has no
intrinsic permanent counterpart: permanents depend on the chosen input and
output port frames, whereas the basis-independent bosonic observable is the
symmetric-cube trace \(h_3(K_T^\dagger K_T)\), a function only of the singular
values.  At each of twenty balanced controls, all six transfers have squared
singular spectrum \(\{1/5,4/5,4/5\}\), giving
\(h_3=313/125\) and exterior-cube probability \(e_3=16/125\); the other
44 Boolean controls have zero fermionic amplitude.  Calibrated determinant
signs form a distance-six simplex code supporting economical three-, five-,
and ten-cut readout, and a rational lossy control realizes an exact chiral
anomaly-free charge vector projectively.  We give a six-mode optical
compilation and quantitative certification thresholds.  Present photonic
components support coherent transfer tomography and bosonic controls; direct
three-fermion emulation additionally requires a totally antisymmetric
three-photon qutrit resource.  The result is a theory and design-limit
analysis, not a report of a built device or a dynamical gauge-theory
simulation.
\end{abstract}

\section{Introduction}

Let \(K\) be a three-input, three-output block of a linear one-particle
transfer.  If all three inputs and outputs are occupied, the fermionic
transition amplitude is \(\det K\), while the collision-free bosonic amplitude
is \(\per K\).  This familiar distinction does not by itself say whether the
two amplitudes encode the same geometric information.  The answer depends on
which changes of input and output coordinates belong to the physical data.

We consider a six-mode transfer obtained from a real symmetric conference
matrix \(C\) of order six, so that
\[
 C^{\mathsf T}=C,\qquad C_{ii}=0,\qquad C_{ij}\in\{\pm1\},
 \qquad C^2=5I.
\]
The two eigenspaces of \(C\), with eigenvalues \(\pm\sqrt5\), both have
dimension three.  A diagonal path control couples them through a
three-by-three block \(K\).  Six coherently marked conference matrices produce
six such blocks.  Their determinants are the six classical Joubert
coordinates and satisfy the two equations of the Segre cubic.  This
identification, including its invariant-theoretic provenance, is developed in
the companion operator paper; here we ask what it means under symmetric and
antisymmetric exchange statistics.

The central point is an orientation boundary.  The Golden data determine the
two three-dimensional port spaces but not ordered orthonormal frames inside
them.  Changing those frames acts on \(K\) from the left and right by
orthogonal matrices.  The determinant retains an orientation sign.  The
permanent does not descend through this action.  Consequently the intrinsic
bosonic companion is not a second signed Segre coordinate: it is a spectral
quantity on the symmetric cube.

Theorem~\ref{thm:exchange-boundary} states this distinction together with the
exact balanced benchmark.  The remaining sections separate intrinsic
observables from calibrated ones, give a short arithmetic application, and
identify the boundary of a photonic implementation.  General
boson--fermion complementarity in linear interferometers is due to
Jabbour and Cerf~\cite{JabbourCerf}; its extension to partially distinguishable
particles is due to Robbio, Jabbour, and Cerf~\cite{RobbioJabbourCerf}.  The
present contribution is the exact Golden specialization, the port-gauge
obstruction, and the associated measurement and resource package.

\section{The Golden transfer}

For one member \(C_T\) of the marked conference family, let
\(Q_{T,+}\) and \(Q_{T,-}\) be matrices whose columns are ordered orthonormal
frames for the \(\pm\sqrt5\)-eigenspaces.  For
\(x=(x_0,\ldots,x_5)\), put
\[
 D_x=\operatorname{diag}(x_0,\ldots,x_5),\qquad
 K_T(x)=Q_{T,-}^{\mathsf T}D_xQ_{T,+}.
\]
When \(|x_i|=1\), the six-mode transfer
\[
 U_T(x)=
 [Q_{T,+}\ Q_{T,-}]^{\mathsf T}D_x[Q_{T,+}\ Q_{T,-}]
\]
is orthogonal.  Values \(|x_i|<1\) describe a diagonal filter followed by
postselection.

The six determinant amplitudes satisfy
\[
 Z_T(x)=\varepsilon_T,10\sqrt5\det K_T(x),
 \qquad \varepsilon_T\in\{\pm1\},
\]
with signs fixed by the marked orientations.  They obey
\[
 \sum_T Z_T(x)=0,\qquad \sum_T Z_T(x)^3=0.
\]
These are the Segre equations in Joubert coordinates.  The equations and the
outer \(S_6\)-action are classical invariant theory
\cite{HowardMillsonSnowdenVakil}; the point used below is their realization as
six amplitudes of one marked transfer.

\section{Why the determinant retains orientation}

Write \(H=K^\dagger K\).  For the intrinsic comparison, define
\[
 B_{\mathrm{sym}}(K)=\operatorname{tr}(\Sym^3 H)=h_3(H),
 \qquad
 F_{\mathrm{ext}}(K)=\operatorname{tr}(\bigwedge^3 H)
 =e_3(H)=\det H.
\]
The first expression is the sum of all-transferred probabilities, averaged
over the choice of normalized three-boson input occupation only after division
by the ten-dimensional symmetric-cube input space.  Thus the operational
uniform-input average is \(B_{\mathrm{sym}}/10\).  The second expression is
the filled three-fermion probability.

\begin{theorem}[Exchange-statistics orientation boundary]
\label{thm:exchange-boundary}
Let \(K_T(x)\) be a Golden transfer block.
\begin{enumerate}[label=\textup{(\roman*)},leftmargin=*]
\item A change of ordered orthonormal port frames sends
\[
 K\longmapsto R_-^{\mathsf T}KR_+,
 \qquad R_\pm\in O(3).
\]
The determinant changes by
\(\det(R_-)\det(R_+)\), while the permanent does not define a scalar on the
resulting double orbit.
\item Every scalar intrinsic under the \(O(3)\times O(3)\) port action factors
through the singular values of \(K\).  Under
\(SO(3)\times SO(3)\), the sign of \(\det K\) is the additional orientation
label.
\item At every balanced control with three entries \(+1\) and three entries
\(-1\), for each of the six marked protocols,
\[
 \operatorname{spec}H=\left\{\frac15,\frac45,\frac45\right\}.
\]
Consequently
\[
 B_{\mathrm{sym}}=\frac{313}{125},\qquad
 F_{\mathrm{ext}}=\frac{16}{125},\qquad
 B_{\mathrm{sym}}-F_{\mathrm{ext}}=\frac{297}{125}.
\]
All intrinsic bosonic scalars are therefore blind to the marked Segre node
and its orientation on this balanced family.
\end{enumerate}
\end{theorem}

\begin{proof}
Part (i) follows directly from changing the two orthonormal frames.  The
determinant transformation law is multiplicative.  The permanent has no
corresponding covariance: for \(K=I_3\), rotating the first two columns by an
angle \(\theta\) changes its value to \(\cos(2\theta)\), although the singular
values and \(|\det K|\) remain fixed.

The real singular-value decomposition classifies the
\(O(3)\times O(3)\) double orbits, proving (ii).  Restricting to the two special
orthogonal groups retains the determinant sign.  For (iii), the Golden
conference identity gives the displayed spectrum on every balanced cut.  The
complete symmetric functions then give
\[
 h_3\left(\frac15,\frac45,\frac45\right)=\frac{313}{125},
 \qquad
 e_3\left(\frac15,\frac45,\frac45\right)=\frac{16}{125}.
\]
Finally, the degree-three symmetric-function identity
\(h_3-e_3=p_1p_2\) gives \(297/125\).
\end{proof}

The theorem does not say that calibrated bosonic measurements lose all
control information.  Once physical port axes are fixed, a permanent is an
ordinary measurable amplitude.  It says that such a value belongs to the
calibrated apparatus rather than to the unmarked Golden operator.

\section{Balanced controls and calibrated readout}

There are \(64\) Boolean controls \(x\in\{\pm1\}^6\).  The twenty balanced
controls are the common nonzero determinant optimum.  For each one,
\[
 |\det K_T|=\frac{4}{5\sqrt5},\qquad
 P_F=|\det K_T|^2=\frac{16}{125}.
\]
The other 44 controls have zero filled-fermion amplitude.  This gives a sharp
Boolean boundary rather than a numerical search criterion.

Freeze the port frames by a reproducible pivot convention.  In that gauge the
collision-free bosonic probabilities on balanced controls take the values
\[
 P_B(111\mid111)\in
 \left\{\frac{16}{3125},\frac{36}{3125},\frac{64}{3125}\right\}.
\]
These values are controls, not invariants.  Coherent access to the calibrated
amplitudes distinguishes all twenty oriented masks; squaring pairs antipodal
controls and loses the overall orientation.

Across the ten projective balanced cuts, the six determinant sign words have
pairwise Hamming distance six.  Three selected signs identify a marked
protocol once Golden membership is already known.  Five signs and magnitudes
jointly test membership and identify the protocol.  Measuring all ten signs
supports soft matched-filter decoding; hard nearest-word decoding corrects
two sign errors.  No relative sign word fixes the global orientation without
a coherent port reference.

The universal complementarity identities for determinants and permanents of
Jabbour and Cerf~\cite{JabbourCerf} provide a useful calibrated consistency
check.  They do not distinguish the Golden family by themselves.

\section{An exact arithmetic instance}

The rational diagonal filter
\[
 x=\left(1,\frac7{13},\frac17,-\frac15,-\frac12,-1\right)
\]
gives
\[
 Z(x)=\frac{72}{455}(11,-10,-8,5,4,-2).
\]
The integer vector on the right has vanishing sum and vanishing sum of cubes,
so it is an anomaly-free chiral charge vector for six Weyl fermions.  General
parametrizations of Abelian anomaly equations are known
\cite{CostaDobrescuFox,CostaDobrescuFoxFew}; the contribution here is the
operator filter and its exact postselection cost.  In the present
normalization,
\[
 P_{F,T}=\frac{(72/455)^2q_T^2}{500},
 \qquad q=(11,-10,-8,5,4,-2).
\]

This construction realizes arithmetic coordinates satisfying the anomaly
equations.  It contains no gauge field, regulator, anomalous current, or
gauge-theory dynamics.  The extension to two \(U(1)\) factors is governed by
lines on the Segre cubic and has its own established algebraic geometry
\cite{GripaiosNguyen}; it is outside the scope of this paper.

\section{Photonic design limit}

The balanced six-mode transfer needs at most fifteen tunable two-mode cells.
A filtered implementation uses two such meshes and six attenuators.  Coherent
transfer characterization fixes the input and output phase gauges and infers
the sign of the reconstructed real determinant.  Six-mode coherent transfer
tomography of this kind has been demonstrated~\cite{RahimiKeshari}, and
programmable twelve-mode processors have operated with three photons
\cite{Somhorst}.

This feasible precursor supports two distinct tests:
\begin{enumerate}[leftmargin=*]
\item reconstruct the six-mode transfer coherently and test the determinant
signs and the linear and cubic Segre residuals;
\item inject ordinary photons and measure the calibrated bosonic controls,
optionally upgrading to the ten-input symmetric-cube average.
\end{enumerate}
The inferred determinant sign in the first test is not a direct measurement
of a many-fermion phase.

The photonic fermion emulator of Matthews et al.~\cite{Matthews} sends an
entangled resource through identical copies of a one-particle process.  For
three particles in this setting it requires a totally antisymmetric state of
three photonic qutrits.  Goyal et al.~\cite{Goyal} give a linear-optical
preparation proposal for this state.  The design evaluated here treats a
demonstrated, characterized source of that resource as an unmet input.  A
three-qutrit GHZ state, although experimentally accessible
\cite{Hu}, has different exchange symmetry and does not substitute for it.

The determinant perturbation bound
\[
 |\det(K+E)-\det K|\leq 3\lVert E\rVert_2
\]
for contraction blocks yields a balanced-sign threshold
\(\lVert E\rVert_2<0.1193\).  The smallest chiral sign requires
\(\lVert E\rVert_2<0.00472\), and ten-percent amplitude precision there
requires \(4.72\times10^{-4}\).  The weak chiral branch, rather than the
balanced benchmark, sets the severe resource requirement.

Diagonal real loss merely changes the real filter and preserves the two Segre
identities.  Coherent diagonal phase error preserves them over
\(\mathbb C\), though a real sign still needs a phase reference.  Transfer
error or mismatch between emulator copies leaves the common carrier.  The
linear residual \(\sum_T\widehat Z_T\) is therefore the first structural null,
followed by the cubic residual \(\sum_T\widehat Z_T^3\).

\section{Discussion}

The Golden transfer separates three levels of information.  Its singular
values determine every intrinsic bosonic scalar.  A calibrated port frame
adds permanent amplitudes and oriented control information.  Antisymmetric
exchange statistics retain the determinant orientation as a one-dimensional
representation of the port-frame action.  The Segre coordinate is therefore
not an arbitrary choice of cubic observable: it is the top exterior amplitude
of the transfer.

The example also separates a feasible optical certification from a direct
many-particle realization.  Current mesh and tomography technology can test
the one-particle carrier and its bosonic shadows.  Direct photonic emulation of
the three-fermion amplitude depends on an additional antisymmetric resource.
This boundary is part of the result rather than an experimental claim.

\section*{Verification}

The exact spectra, balanced census, permanent controls, six-mode compilation,
error thresholds, shot budgets, and decoding schedules are accompanied by
exact generators, compact certificates, SHA-256 manifests, and independent
standard-library replays in the project repository.  The final archival
version will identify a stable artifact release and the precise trust boundary
for each displayed finite claim.

\begingroup
\raggedright
\begin{thebibliography}{99}

\bibitem{JabbourCerf}
M.~G. Jabbour and N.~J. Cerf,
\emph{Boson--fermion complementarity in a linear interferometer},
arXiv:2312.17709.

\bibitem{RobbioJabbourCerf}
M. Robbio, M.~G. Jabbour, and N.~J. Cerf,
\emph{Complementarity between bosonic and fermionic many-body interferences
with partially distinguishable particles},
arXiv:2604.23316.

\bibitem{HowardMillsonSnowdenVakil}
B. Howard, J. Millson, A. Snowden, and R. Vakil,
\emph{A description of the outer automorphism of \(S_6\), and the invariants
of six points in projective space},
J. Combin. Theory Ser. A \textbf{115} (2008), 1296--1303.

\bibitem{CostaDobrescuFox}
D.~B. Costa, B.~A. Dobrescu, and P.~J. Fox,
\emph{General solution to the (U(1)) anomaly equations},
arXiv:1905.13729.

\bibitem{CostaDobrescuFoxFew}
D.~B. Costa, B.~A. Dobrescu, and P.~J. Fox,
\emph{Chiral Abelian gauge theories with few fermions},
arXiv:2001.11991.

\bibitem{GripaiosNguyen}
B. Gripaios and K.~L.~N. Nguyen,
\emph{Anomaly cancellation for two (U(1)) factors},
arXiv:2607.09879.

\bibitem{Matthews}
J.~C.~F. Matthews et al.,
\emph{Observing fermionic statistics with photons in arbitrary processes},
Sci. Rep. \textbf{3} (2013), 1539.

\bibitem{Goyal}
S.~K. Goyal et al.,
\emph{Qudit-teleportation for photons with linear optics},
Sci. Rep. \textbf{4} (2014), 4543.

\bibitem{Hu}
X.-M. Hu et al.,
\emph{Observation of genuine high-dimensional multipartite nonlocality in
entangled photon states},
Nat. Commun. (2025), doi:10.1038/s41467-025-59717-y.

\bibitem{RahimiKeshari}
S. Rahimi-Keshari et al.,
\emph{Direct characterization of linear-optical networks},
Opt. Express \textbf{21} (2013), 13450--13458.

\bibitem{Somhorst}
F.~H.~B. Somhorst et al.,
\emph{Quantum simulation of thermodynamics in an integrated quantum photonic
processor},
Nat. Commun. \textbf{14} (2023), 3895.

\end{thebibliography}
\endgroup

\end{document}
